\begin{abstract}
Persona prompts shape how language models communicate,
but may also affect which requests they refuse.
We evaluate Qwen2.5-7B-Instruct on XSTest under a
neutral baseline and five role personas, producing
2,700 responses. We measure over-refusal of safe
requests and under-refusal of unsafe requests using
an LLM judge, targeted annotation review, and a
separate category for off-topic deflection.
Under the primary labeling scheme, over-refusal
ranged from 0.0\% to 2.8\%, and under-refusal from
5.0\% to 10.5\%. Overall persona differences were
statistically significant, but no individual persona
differed from the neutral baseline after
multiple-testing correction. Conclusions were
sensitive to annotation policy and simulated judge
errors. Our findings suggest that persona prompting
may influence refusal behavior, but provide no
robust evidence that it improves safety calibration.
They also highlight the importance of validating
judgments of persona-dependent response styles.
\end{abstract}
